%! AP Statistics — Unit 3, Lesson 3.7 — Guided Notes KEY
\documentclass[10pt]{article}
\usepackage{apstats-article}
\usepackage{apstats-key}

\begin{document}

\pageheader{Unit 3, Lesson 3.7}{Guided Notes — KEY}
\namedateperiod

\vspace{0.12in}

\begin{objectivebox}
\small By the end of this lesson, I will be able to\dots\\[4pt]
\ans{Define statistical significance and explain what it means in context.}\\[1pt]
\ans{Determine when a causal conclusion is justified.}\\[1pt]
\ans{Use the scope-of-inference table to state what can and cannot be concluded from a study.}
\end{objectivebox}

\vspace{0.1in}

\begin{vocabbox}
\termblanklong{Statistical inference} \ans{Drawing conclusions about a population or treatment based on data collected from a sample or experiment.}
\termblanklong{Statistically significant} \ans{An observed difference so large that it is unlikely to be due to chance alone.}
\termblanklong{Causal conclusion} \ans{Concluding that the treatment \emph{caused} the effect; requires random assignment.}
\termblanklong{Generalization} \ans{Applying results to the broader population; requires random selection.}
\termblanklong{Scope of inference} \ans{The set of valid conclusions (causation and/or generalization) that can be drawn given how the study was designed.}
\end{vocabbox}

\vspace{0.1in}

\begin{notesbox}{Interpreting Experiment Results (VAR-3.E)}
\small
\textbf{Statistical inference:} attributing conclusions from data to the
\ans{population or treatment} from which the data were collected.

\vspace{8pt}
\textbf{Statistical significance:} an observed difference is \textbf{statistically
significant} if it is too \ans{large} to be explained by \ans{chance (random variation)} alone.

\vspace{8pt}
\textbf{Causal conclusion:} If results are statistically significant in an experiment
with random assignment, we can conclude the \ans{treatment} \emph{caused} the effect.
\end{notesbox}

\vspace{0.1in}

\begin{notesbox}{Scope of Inference Table}
\small
\renewcommand{\arraystretch}{2.2}
\begin{tabularx}{\linewidth}{@{} c | X | X @{}}
 & \textbf{Random Selection used} & \textbf{No Random Selection} \\
\hline
\textbf{Random Assignment used} &
Can conclude \ans{causation} \textbf{AND} can \ans{generalize} to the population. &
Can conclude \ans{causation}, but can only generalize to \ans{similar units / the sample}. \\
\hline
\textbf{No Random Assignment} &
Cannot conclude \ans{causation}, but can \ans{generalize} to the population. &
Cannot conclude \ans{causation} and cannot \ans{generalize} beyond the sample. \\
\end{tabularx}

\vspace{8pt}
\textbf{Key rules:}
\begin{itemize}[itemsep=3pt]
  \item Random \textbf{assignment} $\Rightarrow$ \ans{causation}
  \item Random \textbf{selection} $\Rightarrow$ \ans{generalization to the population}
\end{itemize}
\end{notesbox}

\vspace{0.1in}

\begin{practicebox}
\small
\renewcommand{\arraystretch}{1.9}
\begin{tabularx}{\linewidth}{@{} X c c p{2.4cm} p{2.8cm} @{}}
\textbf{Study} & \textbf{Rand. assign.?} & \textbf{Rand. select.?}
  & \textbf{Causation?} & \textbf{Generalize to?} \\
\hline
500 randomly selected adults surveyed; found association between exercise and mood.
& \ans{No} & \ans{Yes} & \ans{No} & \ans{All adults (population)} \\
120 volunteers randomly assigned to meditation or control; significant reduction.
& \ans{Yes} & \ans{No} & \ans{Yes} & \ans{Similar volunteers only} \\
200 randomly selected patients randomly assigned to new drug or placebo; sig.\ improvement.
& \ans{Yes} & \ans{Yes} & \ans{Yes} & \ans{The population of patients} \\
\end{tabularx}
\end{practicebox}

\end{document}
